\documentclass[a4paper,12pt,twoside]{memoir}

% Castellano
\usepackage[spanish,es-tabla]{babel}
\selectlanguage{spanish}
\usepackage[utf8]{inputenc}
\usepackage[T1]{fontenc}
\usepackage{lmodern} % Scalable font
\usepackage{microtype}
\usepackage{placeins}

\RequirePackage{booktabs}
\RequirePackage[table]{xcolor}
\RequirePackage{xtab}
\RequirePackage{multirow}

% Links
\PassOptionsToPackage{hyphens}{url}\usepackage[colorlinks]{hyperref}
\hypersetup{
	allcolors = {red}
}

% Ecuaciones
\usepackage{amsmath}

% Rutas de fichero / paquete
\newcommand{\ruta}[1]{{\sffamily #1}}

% Párrafos
\nonzeroparskip

% Huérfanas y viudas
\widowpenalty100000
\clubpenalty100000

% Imagenes
\usepackage{graphicx}
\newcommand{\imagen}[2]{
	\begin{figure}[!h]
		\centering
		\includegraphics[width=0.9\textwidth]{#1}
		\caption{#2}\label{fig:#1}
	\end{figure}
	\FloatBarrier
}

\newcommand{\imagenflotante}[2]{
	\begin{figure}%[!h]
		\centering
		\includegraphics[width=0.9\textwidth]{#1}
		\caption{#2}\label{fig:#1}
	\end{figure}
}



% El comando \figura nos permite insertar figuras comodamente, y utilizando
% siempre el mismo formato. Los parametros son:
% 1 -> Porcentaje del ancho de página que ocupará la figura (de 0 a 1)
% 2 --> Fichero de la imagen
% 3 --> Texto a pie de imagen
% 4 --> Etiqueta (label) para referencias
% 5 --> Opciones que queramos pasarle al \includegraphics
% 6 --> Opciones de posicionamiento a pasarle a \begin{figure}
\newcommand{\figuraConPosicion}[6]{%
  \setlength{\anchoFloat}{#1\textwidth}%
  \addtolength{\anchoFloat}{-4\fboxsep}%
  \setlength{\anchoFigura}{\anchoFloat}%
  \begin{figure}[#6]
    \begin{center}%
      \Ovalbox{%
        \begin{minipage}{\anchoFloat}%
          \begin{center}%
            \includegraphics[width=\anchoFigura,#5]{#2}%
            \caption{#3}%
            \label{#4}%
          \end{center}%
        \end{minipage}
      }%
    \end{center}%
  \end{figure}%
}

%
% Comando para incluir imágenes en formato apaisado (sin marco).
\newcommand{\figuraApaisadaSinMarco}[5]{%
  \begin{figure}%
    \begin{center}%
    \includegraphics[angle=90,height=#1\textheight,#5]{#2}%
    \caption{#3}%
    \label{#4}%
    \end{center}%
  \end{figure}%
}
% Para las tablas
\newcommand{\otoprule}{\midrule [\heavyrulewidth]}
%
% Nuevo comando para tablas pequeñas (menos de una página).
\newcommand{\tablaSmall}[5]{%
 \begin{table}
  \begin{center}
   \rowcolors {2}{gray!35}{}
   \begin{tabular}{#2}
    \toprule
    #4
    \otoprule
    #5
    \bottomrule
   \end{tabular}
   \caption{#1}
   \label{tabla:#3}
  \end{center}
 \end{table}
}

%
% Nuevo comando para tablas pequeñas (menos de una página).
\newcommand{\tablaSmallSinColores}[5]{%
 \begin{table}[H]
  \begin{center}
   \begin{tabular}{#2}
    \toprule
    #4
    \otoprule
    #5
    \bottomrule
   \end{tabular}
   \caption{#1}
   \label{tabla:#3}
  \end{center}
 \end{table}
}

\newcommand{\tablaApaisadaSmall}[5]{%
\begin{landscape}
  \begin{table}
   \begin{center}
    \rowcolors {2}{gray!35}{}
    \begin{tabular}{#2}
     \toprule
     #4
     \otoprule
     #5
     \bottomrule
    \end{tabular}
    \caption{#1}
    \label{tabla:#3}
   \end{center}
  \end{table}
\end{landscape}
}

%
% Nuevo comando para tablas grandes con cabecera y filas alternas coloreadas en gris.
\newcommand{\tabla}[6]{%
  \begin{center}
    \tablefirsthead{
      \toprule
      #5
      \otoprule
    }
    \tablehead{
      \multicolumn{#3}{l}{\small\sl continúa desde la página anterior}\\
      \toprule
      #5
      \otoprule
    }
    \tabletail{
      \hline
      \multicolumn{#3}{r}{\small\sl continúa en la página siguiente}\\
    }
    \tablelasttail{
      \hline
    }
    \bottomcaption{#1}
    \rowcolors {2}{gray!35}{}
    \begin{xtabular}{#2}
      #6
      \bottomrule
    \end{xtabular}
    \label{tabla:#4}
  \end{center}
}

%
% Nuevo comando para tablas grandes con cabecera.
\newcommand{\tablaSinColores}[6]{%
  \begin{center}
    \tablefirsthead{
      \toprule
      #5
      \otoprule
    }
    \tablehead{
      \multicolumn{#3}{l}{\small\sl continúa desde la página anterior}\\
      \toprule
      #5
      \otoprule
    }
    \tabletail{
      \hline
      \multicolumn{#3}{r}{\small\sl continúa en la página siguiente}\\
    }
    \tablelasttail{
      \hline
    }
    \bottomcaption{#1}
    \begin{xtabular}{#2}
      #6
      \bottomrule
    \end{xtabular}
    \label{tabla:#4}
  \end{center}
}

%
% Nuevo comando para tablas grandes sin cabecera.
\newcommand{\tablaSinCabecera}[5]{%
  \begin{center}
    \tablefirsthead{
      \toprule
    }
    \tablehead{
      \multicolumn{#3}{l}{\small\sl continúa desde la página anterior}\\
      \hline
    }
    \tabletail{
      \hline
      \multicolumn{#3}{r}{\small\sl continúa en la página siguiente}\\
    }
    \tablelasttail{
      \hline
    }
    \bottomcaption{#1}
  \begin{xtabular}{#2}
    #5
   \bottomrule
  \end{xtabular}
  \label{tabla:#4}
  \end{center}
}



\definecolor{cgoLight}{HTML}{EEEEEE}
\definecolor{cgoExtralight}{HTML}{FFFFFF}

%
% Nuevo comando para tablas grandes sin cabecera.
\newcommand{\tablaSinCabeceraConBandas}[5]{%
  \begin{center}
    \tablefirsthead{
      \toprule
    }
    \tablehead{
      \multicolumn{#3}{l}{\small\sl continúa desde la página anterior}\\
      \hline
    }
    \tabletail{
      \hline
      \multicolumn{#3}{r}{\small\sl continúa en la página siguiente}\\
    }
    \tablelasttail{
      \hline
    }
    \bottomcaption{#1}
    \rowcolors[]{1}{cgoExtralight}{cgoLight}

  \begin{xtabular}{#2}
    #5
   \bottomrule
  \end{xtabular}
  \label{tabla:#4}
  \end{center}
}



\graphicspath{ {./img/} }

% Capítulos
\chapterstyle{bianchi}
\newcommand{\capitulo}[2]{
	\setcounter{chapter}{#1}
	\setcounter{section}{0}
	\setcounter{figure}{0}
	\setcounter{table}{0}
	\chapter*{#2}
	\addcontentsline{toc}{chapter}{#2}
	\markboth{#2}{#2}
}

% Apéndices
\renewcommand{\appendixname}{Apéndice}
\renewcommand*\cftappendixname{\appendixname}

\newcommand{\apendice}[1]{
	%\renewcommand{\thechapter}{A}
	\chapter{#1}
}

\renewcommand*\cftappendixname{\appendixname\ }

% Formato de portada
\makeatletter
\usepackage{xcolor}
\newcommand{\tutor}[1]{\def\@tutor{#1}}
\newcommand{\course}[1]{\def\@course{#1}}
\definecolor{cpardoBox}{HTML}{E6E6FF}
\def\maketitle{
  \null
  \thispagestyle{empty}
  % Cabecera ----------------
\noindent\includegraphics[width=\textwidth]{cabecera}\vspace{1cm}%
  \vfill
  % Título proyecto y escudo informática ----------------
  \colorbox{cpardoBox}{%
    \begin{minipage}{.8\textwidth}
      \vspace{.5cm}\Large
      \begin{center}
      \textbf{TFG del Grado en Ingeniería Informática}\vspace{.6cm}\\
      \textbf{\LARGE\@title{}}
      \end{center}
      \vspace{.2cm}
    \end{minipage}

  }%
  \hfill\begin{minipage}{.20\textwidth}
    \includegraphics[width=\textwidth]{escudoInfor}
  \end{minipage}
  \vfill
  % Datos de alumno, curso y tutores ------------------
  \begin{center}%
  {%
    \noindent\LARGE
    Presentado por \@author{}\\ 
    en Universidad de Burgos --- \@date{}\\
    Tutor: \@tutor{}\\
  }%
  \end{center}%
  \null
  \cleardoublepage
  }
\makeatother

\newcommand{\nombre}{Miriam Torres Calvo} %%% cambio de comando
\newcommand{\nombreTutor}{Bruno Baruque Zanón}
\newcommand{\titulo}{Aplicaciones de Visión Artificial en Dispositivos de Edge Computing}

% Datos de portada
\title{\titulo}
\author{\nombre}
\tutor{\nombreTutor}
\date{\today}

\begin{document}

\maketitle


\newpage\null\thispagestyle{empty}\newpage


%%%%%%%%%%%%%%%%%%%%%%%%%%%%%%%%%%%%%%%%%%%%%%%%%%%%%%%%%%%%%%%%%%%%%%%%%%%%%%%%%%%%%%%%
\thispagestyle{empty}


\noindent\includegraphics[width=\textwidth]{cabecera}\vspace{1cm}

\noindent D. \nombreTutor, profesor del departamento de Ingeniería Informática., área de Ciencia de la Computación e Inteligencia Artificial.

\noindent Expone:

\noindent Que la alumna Dª. \nombre, con DNI 45575901K, ha realizado el Trabajo final de Grado en Ingeniería Informática titulado título de TFG. 

\noindent Y que dicho trabajo ha sido realizado por la alumna bajo la dirección del que suscribe, en virtud de lo cual se autoriza su presentación y defensa.

\begin{center} %\large
En Burgos, {\large \today}
\end{center}

\vfill\vfill\vfill

% Author and supervisor
\begin{minipage}{0.45\textwidth}
\begin{flushleft} %\large
Vº. Bº. del Tutor:\\[2cm]
D. \nombreTutor
\end{flushleft}
\end{minipage}
\hfill
% \begin{minipage}{0.45\textwidth}
% \begin{flushleft} %\large
% Vº. Bº. del co-tutor:\\[2cm]
% D. nombre co-tutor
% \end{flushleft}
% \end{minipage}
% \hfill

\vfill

% para casos con solo un tutor comentar lo anterior
% y descomentar lo siguiente
%Vº. Bº. del Tutor:\\[2cm]
%D. nombre tutor


\newpage\null\thispagestyle{empty}\newpage




\frontmatter

% Abstract en castellano
\renewcommand*\abstractname{Resumen}
\begin{abstract}
La tecnologia avanza continuamente y a velocidades que hace unos años eran inexplicables, contamos su ayuda en diferentes entornos de la sociedad y cada vez en mayor medida.
Pero también, contamos con situaciones en las cuáles está no se encuentra tan fácil de acceder, ya que no se cuenta con la facilidad para transportarla y así poder usarla en muchos mas lugares 
y mas comodamente.

Así surge este proyecto, con la idea de poder usar un modelo de Machine Learning en dispositivos de Edge Computing, como puede ser la Jetson Nano de NVIDIA.

Para su desarrollo, se contara con el lenguaje Python y el modelo escogido para entrenar ha sido YOLO en su cuarta versión.
\end{abstract}

\renewcommand*\abstractname{Descriptores}
\begin{abstract}
Deep Learning, Edge Computing, Jetson Nano, YOLO, Python Object Detection, Trasnsfer Learning
\end{abstract}

\clearpage

% Abstract en inglés
\renewcommand*\abstractname{Abstract}
\begin{abstract}
Technology is advancing continuously and at speeds that were inexplicable a few years ago, we count on your help in different environments of society and to a greater extent.
But also, we have situations in which it is not so easy to access, since it is not easy to transport it and thus be able to use it in many more places.
and more comfortably.

This is how this project arose, with the idea of being able to use a Machine Learning model in Edge Computing devices, such as the NVIDIA Jetson Nano.

For its development, the Python language will be used and the model chosen for training has been YOLO in its fourth version.
\end{abstract}

\renewcommand*\abstractname{Keywords}
\begin{abstract}
Deep Learning, Edge Computing, Jetson Nano, YOLO, Python Object Detection, Trasnsfer Learning
\end{abstract}

\clearpage

% Indices
\tableofcontents

\clearpage

\listoffigures

\clearpage

\listoftables
\clearpage

\mainmatter
\capitulo{1}{Introducción}

Descripción del contenido del trabajo y del estrucutra de la memoria y del resto de materiales entregados.

\capitulo{2}{Objetivos del proyecto}

En este apartado se van a presentar los objetivos que han marcado el proyecto, tanto a nivel software, como técnico.

\section{Objetivos Software}
\begin{list}{\textbullet}{ %
    \addtolength{\itemsep}{-2mm} %
    \setlength{\itemindent}{2mm}}

    \item Creación del script que permita entrenar el modelo de detección con las clases seleccionadas, obteniendo como resltado el modelo entrenado (al realizarse mediante YOLO, devolverá un fichero .weights)
    \item Convertir el modelo YOLO (.weights) a un modelo Tensorflow (.pb) para poder trabajar con él.
    \item Creación de los scripts que permitan la detección de los objetos, ya sea partiendo de una imagen o de vídeo (cargando un vídeo o conectando la webcam), los cuáles contarán con diferentes flags de accion durante la ejecución.
    \item Creación de un script de preprocesado de cara a la evaluación de varias imagenes etiquetadas, de tal forma que podamos obtenerla información original en un único fichero.
    \item Medir la calidad de la evaluación del modelo, es decir, calcular el IoU entre las posiciones originales y las predecidas por el modelo, pudiendo obtener su mAP, y resultados sobre la predicción(verdaderos positivos, falsos postivos...), además de que por cada imagen se devolverá la imagen con la posición original, la posición predicha y el IoU.
    \item Mostrar las predicciones en un csv, que muestre el tiempo en el que se ha detectado la predicción, el número de objetos predecidos en dicho instante y el tipo de objeto que es, así como la posición o posiciones en las que se ha encontrado.
    \item Crear un script, que permita identificr las clases con un tracker, es decir, que identifique las clases y nos las vaya etiquetando según vaya detectando.
\end{list}

\section{Objetivos Técnicos}
\begin{list}{\textbullet}{ %
    \addtolength{\itemsep}{-2mm} %
    \setlength{\itemindent}{2mm}}

    \item Convertir el modelo a uno apto para la características de la Jetson Nano.
    \item Usar la plataforma \textit{GitHub} para la organización y gestión del proyecto.
    \item Seguir los principios de la \textit{metodología ágil Scrum}.
    \item Usar herramientas que permitan medir la calidad del código.
\end{list}


\capitulo{3}{Conceptos teóricos}

Para la compresion de este proyecto, se deben conocer los siguientes conceptos:

\section{Deep Learning} 

El deep learning \cite{deepLearning} es una rama del MachineLearning, donde los algoritmos inspirados en el funcionamiento del cerebro humano (redes neurnales) aprenden a partir de 
grandes cantidades de datos y tratan con un alto número de unidades computacionales.

\section{Edge Computing}

El edge computing \cite{edgeComputing} es un tipo de informática que ocurre, en la ubicación fisica del usuario, en la ubicación de la fuente de los datos o cerca de estas. Permitiendo 
que los usuarios obtengan servicios mas rápidos y fiables.

\section{Raspberry Pi}

Una Raspberry Pi \cite{raspberry} es un mini PC de bajo coste, el cual cabe en una mano. Se encuentra compuesto por un SoC, CPU, RAM,puertos de entrada y salida de audio y vídeo, conectividad de red, 
ranura SD para almacenamiento, reloj, una toma para la alimentación, conexiones para periféricos de bajo nivel..., pero a diferencia de un ordenador común no tiene un botón para endecerla/apagarla.

\section{Object Detection}

El Object Detection \cite{objectDetect} es un técnica de vision por ordenador que permite localizar imágenes y/o vídeos. Estos algoritmos se aprovechan del aprendizaje automático o del profundo 
con el objetivo de obtener resultados significativos, es decir, intentan replicar la inteligencia humana a la hora de reconcoer un objeto.
\capitulo{4}{Técnicas y herramientas}

Esta parte de la memoria tiene como objetivo presentar las técnicas metodológicas y las herramientas de desarrollo que se han utilizado para llevar a cabo el proyecto. Si se han estudiado diferentes alternativas de metodologías, herramientas, bibliotecas se puede hacer un resumen de los aspectos más destacados de cada alternativa, incluyendo comparativas entre las distintas opciones y una justificación de las elecciones realizadas. 
No se pretende que este apartado se convierta en un capítulo de un libro dedicado a cada una de las alternativas, sino comentar los aspectos más destacados de cada opción, con un repaso somero a los fundamentos esenciales y referencias bibliográficas para que el lector pueda ampliar su conocimiento sobre el tema.



\capitulo{5}{Aspectos relevantes del desarrollo del proyecto}

En esta sección van a ser detallados los aspectos más relevantes que han tenido lugar a lo largo del desarrollo del proyecto. Al tratarse de un desarrolo \textit{software}, y a lo largo de este se han encontrado diferentes retos e inconvenientes. A su vez, se han aplicado buenas prácticas con el objetivo de obtener un resultado de calidad.

\section{Investigación}
El eje principal que posee el proyecto es el \textit{Object Detection}, el cuál a pesar de ser conocido, nunca se había trabajado con él, se puede decir que su conociemnto era puramente teórico. En la misma línea, se desconcoia completamente el funcionamiento de la \textit{Jetson Nano}, la cuál a trabajar de forma intenra con \textit{Ubuntu} ha supuesto
una gran ayuda y beneficio al ser un sistema operativo previamente conocido.
Por todo esto, ha supuesto un doble esfuerzo, especialmente por el primero de ellos, ya que se ha requerido una formación previa par poder conocer el correcto funcionamiento a la hora de entrenar un modelo de detección, así como averiguar cuál era la mejor opción de cara al proyecto.
De cara al objetivo de obtener un mator conocimento acerca de \textit{Object Detection} se leyeron diferentes \textit{papers} \cite{Pathak2018,Yang2017, Wu2020}, obteniendo de ellos los suficientes conocimientos sobre su funcionamiento, la calidad de la detección y el redimiento.

\section{Metodología \textit{Scrum}}
Tal y coo se comento en el punto \ref{scrum}, el proyecto se ha reliazado siguiendo una metodología ágil, lo cuál nos permite trabajar con \textit{sprints}, de tal manera, que el trabajo a relaizar en cada uno de ellos, se encuentre documentado desde el inicio, para así poder con una mayor eficiencia, pudiendo priorizar las tareas en función de las existentes y contando con diferentes versiones según a la vez que se va siguiendo el desarrollo del proyecto.

Los conocimeientos que se poseían de \textit{Scrum} eran más bien teóricos, pero si que se había tenido la oportunidad de trabajar con ella, durante la realización de las prácticas en empresa.

Una de las principales dificultades encontradas, ha sido la estimación del tiempo que va a ser necesario para pdoer llevarlo a cabo, todo ello apartir de un nombre y descripción, lo cuál ha concluido con resultados no muy exactos, ya que se en algunos casos la estiamción se ha realizado por debajo y en otros por arriba. Si a esto se le suma que ciertas tareas han llevado más tiempo del debido a problemas versiones y de compatibilidades entre sistemas, hay que suamrle que al contar con un determiando número de asiganturas
más la realización de manera paralela de las prácticas, ha conllevado que tareas que de normal llevan pcoo tiempo, han llevado más de lo planeado.

\section{Creación del dataset de entrenamiento y entremiento del modelo de detección}
A lo largo del Grado, no se ha visto nada relacionado con la \textbf{Inteligencia Artificial}, siendo esto una desventaja a la hora de crear el dataset del modelo y etiquetarlo con las características correspondientes al algoritmo de detección escogido. Lo que ha conllevado una inversión de horas, de cara a la búsqueda de infromación, así como de las diferentes formas para entrenar el modelo, y con ello la más asequible de cara a requisitos técnicos, debido a que llevar a cabo un entrenamiento de estas características, 
requiere de una gráfica potente, así como un equipo los suficientemente potente para llevarlo a cabo, ya que el modelo necesita estar un mínimo de horas entrenando para que el resultado sea medianamente decente y ser funcional.
Debido a esto, las labores de entremiento, se han realizado sobre \textbf{Google Colab}, visto en el Apartado \ref{colab}.

\section{Pruebas realizadas}
A la hora de probar la funcionalidad de modelo que se ha entrenado, se han realizado diferentes pruebas de detección tras su conversión a un modelo \textit{.pb}, el cuál se basa en el grafo de inferencia que calcula \textit{Tensorflow}.
Seguidamente, se comprobó mediante ejecuciones que el modelo es funcional, es decir, que detecta las clases, para las cuáles ha entrenado, detectando imágens y vídeos.
Con rwsultados como los que se ven a continuación:

\imagen{coco_detection}{Detección a través del modelo COCO}
\imagen{license_plate_detection}{Detección a través del modelo creado para detectar matrículas}
\imagen{american_plate_detection}{Detección a través del modelo creado para detectar matrículas 2}
\imagen{head_detection}{Detección a través del modelo creado para detectar cabezas}

Tras esto, se realizaron scripts para medir el \textit{accuracy} y para preprocesar las posiciones originales, es decir, castearlas y devolverlas en un único fichero que contenga la ruta de la imagen, el parámetro \textit{x1}, el parámetro \textit{y1}, el parámetro \textit{x2}, el parámetro \textit{y2} y por último el id de la clase que se ha detectado en dicha posición.
Además, por cada imagen que detecta la devuelve representando la posición original, la posición representada por el modelo y el \textit{IoU} entre ambas posiciones.

\clearpage

Tal y como se muestra a continuación junto con los resultados de la evaluación
\imagen{iou2}{Representacioón de la posición original\textit{(verde)}, la posición predecida\textit{(azul)} y el valor del IoU\textit{(amarillo)}.}
Resultados de la evaluación de la calidad del código con un IoU de 0,70:
\imagen{mAP_lp}{mAP del modelo detector de matrículas} 
\imagen{license_plate_graph}{Gráfico \textit{Precission} y \textit{Recall} de la clase \textit{license plate}}
\imagen{info_license_plate_predict}{Información sobre los \textit{False Predictions} y \textit{True Positives}}

% \section{Jetson Nano}

\section{Resultados evaluación de código} 
Durante el desarrollo se integró una herramienta que permitiesen medir la calidad del código, con el objetivo de controlar la calidad del desarrollo, respecto a los estándares definidos (\textit{ver Apartado }\ref{calidad_codigo}).
En este aprtado veremos los resultados que se han ido obteniendo, los cuáles se ejecutaban automáticamente los análisis al realizarse un \textit{commit}.

\imagen{code_quality_results}{Resultados de la evaluación de la calidad del código}

\clearpage

A través, del dashboard que posee \textit{SonarCloud} para el proyetco, podemos obtener más información sobre la calidad y vulnerabilidades que posee el código.

\imagen{dashboard_sonar}{Resultados mostrados desde el \textit{Dashboard} de SonarCloud}

\imagen{evaluacion_sonar}{Gráfico que muestra la evolución del código}

Gracias al uso de esta herramienta, se ha conseguido reducir los \textit{code smells} hasta el punto de poseer una calificación de A.
\capitulo{6}{Trabajos relacionados}

Este apartado sería parecido a un estado del arte de una tesis o tesina. En un trabajo final grado no parece obligada su presencia, aunque se puede dejar a juicio del tutor el incluir un pequeño resumen comentado de los trabajos y proyectos ya realizados en el campo del proyecto en curso. 

\capitulo{7}{Conclusiones y Líneas de trabajo futuras}

En este apartado se van a agrupar tanto las conclusiones a las que se llega tras el desarrollo del proyecto, como los posibles puntos de mejora, como las líneas de trabajo futuras.

\section{Conclusiones}
Conclusiones a las que se llega tras el posterior desarrollo del proyecto.
\begin{itemize}
    \item Los objetivos del proyecto se han cumplido satisfactoriamente
    \begin{list}{\textbullet}{ %
        \addtolength{\itemsep}{-2mm} %
        \setlength{\itemindent}{2mm}}
        \item \texttt{Object Detection}. Posee un algorimto para poder entrenar cualquier modelo con el algoritmo YOLO en la versión 4. Así como, un script para convertir dicho modelo a uno apto para Tensorflow, y así poder detectar vídeos, imágenes y/o trabajar con Obejct Tracking.
        A su vez, podemos exportar las posiciones de cada uno de estos scripts a un fichero \textit{csv}, como guardar el resultado deseado en un vídeo o en una imagen.
        Pero sobre todo, nos permite medir la caldiad del modelo entrenado, con el objetivo de evaluar su calidad de evalaución.  
    \end{list}
    \item La lectura de artículos científicos ha sido algo prácticamente nuevo, siendo necesarias numerosas horas para su interpretación, se debe destacar que según se avanzaba en el desarrollo del proyecto, la familiarización con éstos ha sido satisfactoria, permitiendo asimilar la información como la idea que se iba a implementar.
    \item El trabajar con Object Detection, a través del algoritmo de detección de YOLO, permite que sea rápido y su calidad de detección sea bastante buena.
    \item La creación de un modelo de detección, así como su respectivo dataset de entremiento ha permitido aumentar el conocmiento que se tenia sobre la detección de objetos, así cmo trabajar con herramientas que al inicio del proyecto eran totalmente desconcoidas. 
\end{itemize}

Como \textbf{conclusión} cabe destacar, la oportunidad tanto de aprendizaje como la de trabajar en un proyecto nuevo ha supuesto. Desde el inicio ha sido necesario dominar ciertas tecnologías vistas (y nuevas) a lo largo de todo el Grado, incluso a la hora de escribir la documentación.

\section{Líneas de trabajo futuras}
Dadas las características del proyecto, existen diferentes líneas de mejora y trabajo futuras.

\begin{list}{\textbullet}{ %
    \addtolength{\itemsep}{-2mm} %
    \setlength{\itemindent}{2mm}}
    \item Crear scripts de entrenamiento, conversión y adaptación de los scripts de detección a YOLOv5.
    \item Permitir la importación y conversión a otros algoritmos diferentes de YOLO y sus versiones.
    \item Crear funciones que permitan apoyar la detección de los objetos, y en función de la clase o clases que se vayan a detectar se realize unas u otros, por ejemplo, en el detector de matrículas que permita leerlas y exportalas a un fichero, en modelos que permitan detectar vehículos contar vehículos en ambas direcciones, con el objetivo de contar los que entran y salen de una carretera y exportarlos a un fichero \textit{csv}.
    \item Realizar un proceso de refactorización código, ya que si el proyecto crece y por ende, aumentan sus funcionalidades.
    \item Tras la adaptaciones a otros algoritmos de detección, adaptarlo en función de cada algortimo, con el objetivo de tratarlo como una librería y subirlo a PIP \cite{pip}.
\end{list}


\bibliographystyle{IEEEtran}
\bibliography{bibliografia}

\end{document}
